\section{就业、开源与行业格局}

\subsection{翻转比特比加速物质快一百万倍}

Karpathy 提出了一个分析框架：当前正在发展的主要是\textbf{数字 AI}——可以在数字世界中操纵信息的"幽灵实体"，目前没有物理实体。翻转比特和复制粘贴数字信息，比加速物质快一百万倍。

因此他预测了一个三阶段路线图：
\begin{enumerate}
    \item \textbf{数字空间大规模解除束缚（unhobbling）}——之前因为人类思考周期不够而没被充分处理的数字信息，会被 AI 大量重写
    \item \textbf{数字与物理的接口}——传感器（看见世界）和执行器（对世界做点什么）。很多有趣的公司会出现在这个接口上
    \item \textbf{物理世界的全面渗透}——市场规模可能远大于数字空间，但难度也大一百万倍
\end{enumerate}

\begin{knowledgebox}{杰文斯悖论（Jevons Paradox）与软件需求}
经典案例：当年很多人担心 ATM 会取代银行柜员，但 ATM 降低了银行网点的运营成本，于是开了更多网点，反而雇了更多柜员。Karpathy 认为类似的事情正在软件领域发生：软件变便宜了，海量被压抑的需求会释放。代码现在是"临时性的"、可修改的——整个数字基础设施有巨大的重写需求。
\end{knowledgebox}

他对软件工程"谨慎乐观"——需求弹性等因素会起作用，但也坦承自己在预测方面不是专业的。

\begin{quote}
``I went around OpenAI and I was like, you guys realize if we're successful, we're all out of a job. We're just building automation for Sam or something like that.''

——我在 OpenAI 里到处跟同事说，你们知道吗，如果我们成功了，我们所有人都失业了。我们本质上就是在给 Sam 或者董事会造自动化系统。
\end{quote}

\subsection{开源：意外地落在了一个还不错的位置}

开源与闭源前沿的差距在持续缩小：从最初的完全没有，到落后18个月，到现在大约落后6到8个月。

Karpathy 用 Linux 做类比——Linux 运行在绝大多数的服务器上，因为行业需要一个共同的开放平台。AI 领域也是同样的需求。两个挑战：
\begin{itemize}
    \item 资本支出是硬约束——训练前沿模型需要大量资金
    \item 前沿智能的需求依然存在——诺贝尔奖级别的工作或巨型项目可能是闭源实验室的地盘
\end{itemize}

\begin{quote}
``Centralization has a very poor track record in my view.''

——中心化的历史记录在我看来很差。
\end{quote}

他提到东欧的历史教训，以及机器学习中的集成模型（ensemble）永远优于任何单一模型。他希望当最难的决策到来时，有更多人在房间里。

\begin{warningbox}{进一步中心化的隐忧}
Karpathy 指出了一个让他不安的趋势：即使在闭源端，领先者的圈子也在进一步缩小。不是所有前沿实验室都保持在最顶尖。他希望有更多实验室，有更多声音参与决策。
\end{warningbox}

\subsection{为什么不回前沿实验室}

Sarah 替 Noam Brown 问了一个问题：你明明可以在前沿实验室里做 AutoResearch，有大规模算力，有同事——为什么不去？

Karpathy 列出了几个原因：

\textbf{第一，利益冲突。}你在前沿实验室有巨大的财务激励。而你自己承认 AI 将以极其戏剧性的方式改变人类社会。你在这里一边造这项技术、一边从中获益、一边被财务手段深度绑定——这个矛盾正是 OpenAI 最初试图解决的问题，至今没有真正解决。

\textbf{第二，你不是完全自由的主体。}

\begin{quote}
``If you're inside one of the frontier labs, there are certain things that you can't say, and conversely there are certain things that the organization wants you to say. They're not going to twist your arm but you feel the pressure.''

——在前沿实验室内部，有些话你不能说，有些话组织希望你说。不会扭你的胳膊，但你能感受到那种压力。
\end{quote}

\textbf{第三，员工对组织行为的实际影响力有限。}当真正的高风险时刻到来时，你在房间里贡献想法，但你并不真正掌控那个实体。

他认为理想方案可能是阶段性进出：进去一段时间做有价值的工作、了解实际进展，然后出来保持独立性。他也承认离开后判断力会漂移——实验室是不透明的，你对系统实际工作原理的理解会逐渐过时。

\subsection{本章小结}

数字空间将率先经历大规模改造，物理世界滞后。软件工程可能因杰文斯悖论而需求增长。开源落后闭源6--8个月是"还不错的位置"。前沿实验室的中心化趋势令人担忧，Karpathy 选择留在外面是因为独立性和利益冲突的考量。


